\section{Everything Solid Melts into Air}

\begin{quotex}
All fixed, fast-frozen relations, with their train of ancient and venerable prejudices and opinions, are swept away, all new-formed ones become antiquated before they can ossify. All that is solid melts into air, all that is holy is profaned.
\flright{\textsc{Karl Marx}, \emph{Communist Manifesto}}
\end{quotex}


In these brief yet portentous sentences, Marx states the essence of the spirit of communism that haunts the West. All values and social relations that a well-bred man would have considered to be sane and normal are now denied. Yet even as new relations arise to displace the old, they in turn will give way to new ones before they can ossify. Thus, the progressive is never satisfied, there is always some novelty just over the horizon. No progressive held it against Obama, for example, when he opposed same sex marriage. As expected, Obama ``evolved'' on the issue, without explanation, without accusations of ``flip flopping''.

Everything solid refers to formal causes, the notion that people, conditions, things, and events are manifestations of fixed ideas. The ideas evaporate and take abode in some non-existent heaven, leaving their former manifestations without firm boundaries, and thus merge into each other in an ever increasing lack of manifestation. Of course, all that is holy also evaporates; however, it is insufficient to let it run its natural course, rather it must be mocked, profaned, and combated.

At one time I knew many Marxists, but today no one admits to being one. However, I know how and what Marxists think, and many think the same thoughts today without realizing their true source. The genius of Marx is that he understood how this process would take root, not in undeveloped areas as Lenin and Mao presumed, but rather as the natural development of advanced capitalist societies.

\paragraph{The Spectre of Communism}


Marx perceived clearly the spectre of communism, or we would say its egregore\footnote{\url{https://en.wikipedia.org/wiki/Egregore}}; he nourished it, defined it, promoted it, until subtly it has grown in psychic strength and manifested as the wandering influence in the minds of many in our day. An egregore is a psychic thoughtform infecting a like-minded group of human beings from whom it draws its life. It then amplifies the thoughts of its component members, returning to them much more strongly. That is why the thoughts of progressives have such a high affective element, while lacking intellectual coherence.

The egregore controls their thoughts, which they are unaware of. That is why progressives can suddenly, and even simultaneously, adopt the same cause without outward collusion. As the Proverb says, ``the locusts fly in formation even though they have no king.'' And the devouring locusts are relentless, sweeping aside indiscriminately every venerable and hoary truth and profaning anything holy.

\paragraph{Social and Political Disorder}


Now consider these general principles. Are they held only by open communists, or to some extent by anyone who calls himself a liberal or a progressive? The answer is clear. Of course, by the logic of progressivism, those who balk today at the more extreme elements of these principles may very well adopt them tomorrow.

\begin{itemize}


\item They support every revolutionary movement against the existing social and political order of things.

\item They disdain to conceal their views and aims. They openly declare that their ends can be attained only by the forcible overthrow of all existing social conditions.

\item They attack every principle of existing society.

\item Abolition of the family
\end{itemize}


The progressive intellectuals attack the bourgeoisie from above while supporting every movement that attacks it from below. In the USA, there are increasingly loud voices advocating forcible overthrow: eliminate the constitution, ignore the bill of rights, rule by decree (executive order or regulations) rather than law, appeal to the masses directly over congress, destroy political opponents.

What principle is immune? Even the boy scouts must be coerced to yield. The family is attacked by calling everything ``family''. Existing (bourgeois) society likes guns, since it has always been the right of a free man to bear arms. Hence, guns must be opposed. Existing society is completely dependent on petroleum energy, hence oil must be opposed. I'm sure readers can easily add to this list.

\paragraph{The Battle against the Bourgeoisie}


The battle of the progressive and communist left is against the bourgeoisie because it is the source of the moral and physical strength of previously existing society. Since it was European men who created this bourgeois society, they are deemed to be the enemy. This brings in a racial element, since it seems that white men are the enemy, but that element is derivative, an epiphenomenon. White progressives are not considered to be the enemy, just those who persist in maintaining the existing order of things. To make race the primary element is to misunderstand the essence of the fight; that will lead to wrong-headed efforts to oppose the revolution.

\paragraph{The Battle against Christianity}


The bourgeoisie is also Christian, or at least culturally Christian. Hence, the revolution must be openly anti-Christian. The evidence for this is abundant, even among those who are not overtly or consciously Marxist. Unfortunately, those who only see the racial component are often themselves anti-Christian, thus inadvertently aiding the revolution.

This battle is not only open, but also subversive. Marx cynically claimed:

\begin{quotex}
Nothing is easier than to give Christian asceticism a Socialist tinge.
\end{quotex}


In this way, the revolution can recruit, and has recruited, Christian elements to do this work. Of course, in Tradition, Christian asceticism applies only to those who choose such a lifestyle, typically as monks. The counsels of perfection include obedience, poverty and celibacy. Clearly socialism is not promoting any of those counsels, so this alleged socialist tinge is bogus, albeit very effective.

\paragraph{The Bourgeois Counterrevolution}


Since the power of the bourgeoisie was the result of a revolution against the earlier feudal system, it cannot act as the bulwark against the revolution. That is why bourgeois parties always continue the leftward drift. In effect, they already contain some of the premises of Marxism. We see why Marx was led to conclude that Marxism follows inevitably from capitalism, and not from the earlier spiritual tradition which was opposed to both revolutionary movements. We can list three here as examples.

\subsubsection{Internationalism}


Bourgeois capitalism is thoroughly international, and Marx mentions this several times in the Manifesto, if anyone cares to look it up.

\subsubsection{Economism}


For bourgeois capitalism, the economy is primary. It even serves as the highest unifying principle, transcending even spiritual matters. Obviously, this is contrary to Tradition, for which the economic class is subservient to the spiritual authorities. We can refer to Voltaire in this matter:

\begin{quotex} 
At the London Stock Exchange, men of all religions treat with one another without asking in whom or in what they believe and give the name of infidel only to those who go bankrupt.
\end{quotex}


\subsubsection{Women in Common}


One of the more scandalous ideas of the Manifesto was that of women in common. Obviously, that has come to pass and now appears normal. But, that was not so in the 19\textsuperscript{th} century. Yet Marx shrewdly observed:

\begin{quotex}
Bourgeois marriage is, in reality, a system of wives in common.
\end{quotex}


Divorce and remarriage is an example that system, and is widely accepted even among the bourgeois elements today.

\paragraph{Economic Untenability}


The bourgeois parties of today frame their arguments in terms of the untenability of the programs of the left. They naively argue that the national debt cannot increase without limit, that green energy programs will not work, that gun control will not reduce crime, and so on. All this is true, but sadly misses the point. The left claims to be based on science and intelligence, which point to progress; hence, anything that opposes that progression is \emph{ipso facto} unscientific. In his recent inauguration address, Obama outlined a program, very little of which could be justified on rational and objective scientific grounds. Marx anticipated this, and the locusts follow, consciously or unconsciously:

\begin{quotex}
Of course, in the beginning, this cannot be effected except by means of \emph{despotic inroads on the rights of property}, and on the conditions of bourgeois production; \emph{by means of measures}, therefore, \emph{which appear economically insufficient and untenable}, but which, in the course of the movement, outstrip themselves, necessitate further inroads upon the old social order, and are unavoidable as a means of entirely revolutionising the mode of production.\footnote{My emphasis.}
\end{quotex}


The measures are not intended to be sufficient and tenable; hence, the bourgeois party misses the point totally and fails as a rational opposition. Programs such as progressive income tax, government-controlled healthcare, green energy, carbon credits, food stamps, and so many others are not intended to promote free and prosperous economic activity. Their intent is simply to make ``inroads'' into free activity, to limit such activity, with the ultimate goal of ``revolutionising the mode of production''. This is the ``open secret'' of Marxism.


\flrightit{Posted on 2013-01-29 by Cologero}

\begin{center}* * *\end{center}

\begin{footnotesize}
\begin{sffamily}
\texttt{Logres on 2013-01-29 at 10:55 said:}

Ugh. I wish I could deny what you say here, but it is all too perspicacious. The insight on race was helpful -- many blacks where I live are by nature highly (and truly) conservative, much more so than most whites (in general). The same is true for Mexicans. They are simply being herded where the managerial elites and their egregore want them to be. It's useful to stir up more strife, as it tends to give everyone an apocalyptic feeling that makes them more subject to manipulation. I've always wondered why the rules don't apply to the Left (even when it is their own rules) and this article makes it crystal clear : there is merely affective hate for existing orders, so ``you have to break some eggs to make omelets''. Nothing to put in its place.

\hfill

\texttt{Janus on 2013-01-29 at 17:31 said:}

I have to say, the argument that ``liberals and progressives'' hold similar views among those you outlined as communist and may draw more radical conclusions tomorrow feels like a bit of a cop-out. It might be said of Franco, ``he to the extent that he upheld National Syndicalism was in favor of worker control of businesses''. That sounds communist. But whatever else, he was certainly no communist. Abolition of the family would not happen under him.

This article does seem to leave out the role of the bourgeoisie in the steps toward communism, particularly that it was they responsible for the revolutions which overthrew major institutions of the Traditional order, such as monarchies and the Church. This is the climate within which Marx was writing. It is they too who are responsible for the reversal of those remaining values from a `modern' Christian form to the post-60's individualism, false liberation and the overthrow of spiritual authority even from the realm of values. Similarly, a Traditionalist can support ``environmentalist'' causes like sustainable energy, a reigning in of the capitalist titans, and a better role for workers in an economy without being a Stalinist in the making (think along Ezra Pound lines here). It is an unfortunate habit of many Traditionalists that we will sometimes defend aristocrats and priests who ``exceed their brief'', and then attack the workers and merchants who react against this. Think of the patricians of the Roman Republic who used their power to get hold of lands which were supposed to belong to the people of Rome. They then accused Gracchus, who wanted this land redistributed among citizens, of being the one who was destroying the Republic. I'm not saying you're doing this here, but it's a tendency we must guard against.

\hfill

\texttt{Cologero on 2013-01-29 at 18:36 said:}

Thanks for your comments, Janus, but perhaps we need to read understand the argument more carefully, especially in the context of everything that has been written previously. The example you chose from Franco is on the level of particular policies, while we were addressing the fundamental thought processes in play, even if not explicitly understood by their adherents.

We have mentioned many times, following Evola's lead, that the mass of men do not think in an active way, but rather thoughts pass through them. We have mentioned many times, following Guenon, that there are psychic forces in play, what he sometimes called ``wandering influences'', that influence men without their conscious awareness; he even wrote about being under the unconscious influence of Satan (not to take a overly lurid interpretation of that).

Read the article again, please. There was actually a section devoted to the idea that the power of the bourgeoisie was the result of a revolution; we pointed out that it makes them ineffective as a counter-revolutionary force. We pointed out that the presuppositions of bourgeois thought are similar to the marxists, so they cannot provide an effective resistance to progressive movements.

I don't recall attacking workers and merchants. Although I didn't specifically address in this article a Traditional economy, I did allude to free activity. I assumed everyone would have understood that as free from unnecessary autocratic control.

Please try to understand our program, Janus. We are not interested per se in particular historical events, insofar as principles cannot be derived from them. As we wrote recently in ``Thought Power'', it is a question about exposing the deep and fundamental presuppositions of thought. These presuppositions are usually invisible, even to those who hold them. Only the few great thinkers think at that level. Marx is one of them, along with some others we have dealt with, e.g., Bacon and Hobbes. They have created the modern world by consciously rejecting certain Traditional principles; the mass of people absorb those forms of thought unconsciously. Why this is so is a complex topic.

Janus, one of the rules is that explicit references to the text must be made. So if we have an ``unfortunate habit'', or unguarded tendencies, please show me exactly where you see that; of course, a habit implies it happened more than once.

In light of this, would you like to revise your comment?

\hfill

\texttt{Jason-Adam on 2013-01-29 at 22:49 said:}

By the definitions you use above, Cologero, I propose that Alexandr Dugin, Alain de Benoist and the other ``new right'' ``fourth political theorists'' ``global revolutionaries'' are Communists in disguise.

One thing I am very grateful for is how in your translations of Evola we get to see just how deeply conservative he was, whereas the Duginists are trying to remake him into some kinda leftist figure.

\hfill

\texttt{Janus on 2013-01-30 at 03:02 said:}

You did, in fact, reference the bourgeois revolution. Apologies, I didn't read that carefully enough. In terms of referencing the text, I'll just point to what you say about the egregore in the beginning, that it is a psychic projection of sorts. Is this then a sort of counter-form of the Idea, rising from below rather than coming from a higher plane of reality? If so, would you say that these have a sort of consciousness and life of their own (and could thus be called demons) or that these represent unconscious tendencies in the psyche, an inversion of an Idea which people do not realize they are being influenced by and which people could potentially become aware of and thus exorcise (which, thinking about it, seems quite demonic as well)? Insofar as we are focusing on principles, I think this suggests that one must be weary that actions which might be just (say, fighting against violations of liberty) are not influenced by false principles, and thus perhaps by an unknown influence, which the demonic egregore appears to be. I am uncertain form this text whether the egregore is limited to negative influences or can be positive as well (a ``good'' egregore... or would this be the Idea?).

\hfill

\texttt{Janus on 2013-01-30 at 03:06 said:}

From my understanding of the article, for this one would have to prove hat Dugin and de Benoist are indeed reaching their conclusions under the influence of the communist egregore. Seeing as Dugin does claim to take Tradition as a foundation point, I would say it would be best to give him some benefit of the doubt and first examine his principles and what he concludes from them before charging him with atheistic communism. He could even be reaching some faulty conclusions from valid and true principles, in which case education is far more important. I must say, I have never seen any new Rightist or 4PT advocate claim that Evola was in any way a leftist. Can you give a reference?

\hfill

\texttt{Cologero on 2013-01-31 at 21:51 said:}

Everything casts a shadow; the egregore, too, casts a shadow. It is not so easy to find the light in a world of shadows. Like the song of the sirens, life in the shadows is tempting, it promises so much but delivers so little. To remain standing, to stay conscious, is to answer all such questions.


\hfill

\end{sffamily}
\end{footnotesize}

